\pagestyle{plain}

This section introduces the notion of Jordan recoverability on subcategories of representations of quivers, firstly studied by \cite{GPT19}. In the following, we will not use examples to illustrate the concepts. We refer the reader to \cite{GPT19,D22,D23} for deeper details.

\begin{conv} \label{conv:addsubcat}
    Our subcategories are full and closed under sums and summands
\end{conv}

\subsection{Quiver representations}

A \new{quiver} is a quadruple $Q = (Q_0,Q_1,s,t)$ where:
\begin{enumerate}[label=$\bullet$, itemsep=1mm]
    \item $Q_0$ is a set called \emph{vertex set};
    \item $Q_1$ is a set called \emph{arrow set}; and,
    \item $s,t : Q_0 \longrightarrow Q_1$ are functions called \emph{source} and \emph{target functions}.
\end{enumerate}
Such a quiver $Q$ is said to be \new{finite} whenever $Q_0$ and $Q_1$ are finite sets. In the following, we assume that all our quivers are finite.

A \new{path} of $Q$ is either a formal variable $e_q$ for some $q \in Q_0$ called \emph{lazy path at $q$}, or a finite nonempty word $p = \alpha_k \cdots \alpha_1$ on $Q_0$ such that, for all $i \in \{1,\ldots,k-1\}$, we have $s(\alpha_{i+1}) = t(\alpha_i)$. Given such a path $p$, we define:
\begin{enumerate}[label=$\bullet$, itemsep=1mm]
    \item its \emph{source} $s(p)$ to be $s(\alpha)$ (we set $s(e_q) = q$);
    \item its \emph{target} $t(p)$ to be $t(\alpha_k)$ (we set $t(e_q) = q$); and,
    \item its length $\ell(p)$ to be $k$ (we set $\ell(e_q) = 0$).
\end{enumerate}
For any arrow $\alpha \in Q_1$, we consider its formal inverse $\alpha^{-1}$ such that $s(\alpha^{-1}) = t(\alpha)$ and $t(\alpha^{-1}) = s(\alpha)$. We denote by $\overline{Q_1}$ the set of formal inverses of arrows in $Q_1$. A \new{walk} of $Q$ is either a lazy path or a finite nonempty word $w = \beta_k \cdots \beta_1$ on $Q_1 \sqcup \overline{Q_1}$ such that, for all $i \in \{1,\ldots,k-1\}$, we have $s(\beta_{i+1}) = t(\beta_i)$. We extend the source, target, and length functions to walks of $Q$ similarly to the paths of $Q$. A quiver $Q$ is said to be \new{connected} whenever for any pair $(a,b) \in Q_0$, there exists a walk $w$ of $Q$ such that $s(w) = a$ and $t(w)=b$. From now on, we assume that all our quivers are connected.

Let $\mathbb{K}$ be a field. For geometric purposes, we assume that $\mathbb{K}$ is algebraically closed. The \new{path algebra} of $Q$, denoted by $\mathbb{K}Q$, is the $\mathbb{K}$-vector space generated as a basis by the set of all the paths of $Q$, together with a multiplication which acts as concatenation on the paths: meaning \[p_2 \cdot p_1 = \begin{cases}
    p_2p_1 & \text{if } s(p_2) = t(p_2) \text{; and,} \\
    0 & \text{otherwise.}
\end{cases}\] For any $m \in \mathbb{N}$, we denote by $\mathbb{K}Q_{\geqslant m}$ the (bi-sided) ideal of $\mathbb{K}Q$ generated by all the paths of length greater than or equal to $m$. An ideal $I \subseteq \mathbb{K}Q$ is said to be \new{admissible} whenever there exists an integer $m \geqslant 2$ such that $\mathbb{K}Q_{\geqslant m} \subseteq  I \subseteq \mathbb{K} Q_{\geqslant 2}$. A \new{relation} on $Q$ is a (finite) $\mathbb{K}$-linear combination of paths of length at least two, and having a common source and a common sink. For any subset of relations $R$ on $Q$, we denote by $\langle R \rangle$ the ideal of $\mathbb{K}Q$ generated by $R$. Call \new{bounded quiver} any pair $(Q,R)$ where $Q$ is a quiver and $R$ is a set of relations on $Q$.

Given a bounded quiver $(Q,R)$, we define a \new{representation} of $Q$ (over $\mathbb{K}$) as a pair $E=((E_q)_{q \in Q_0}, (E_\alpha)_{\alpha \in Q_1})$ such that:
\begin{enumerate}[label=$\bullet$, itemsep=1mm]
    \item $E_q$ is a $\mathbb{K}$-vector space, for all $q \in Q_0$;
    \item  $E_\alpha : E_{s(\alpha)} \longrightarrow E_{t(\alpha)}$ is a $\mathbb{K}$-linear transformation, for any $\alpha \in Q_1$; and,
    \item by setting $E_p = E_{\alpha_k} \circ \cdots \circ E_{\alpha_1}$ for any path $p = \alpha_k \ldots \alpha_1$ of $Q$, whenever we have a relation $k_1 p_1 + \ldots + k_s p_s$ on $Q$ in $R$, then $k_1 E_{p_1} + \ldots  + k_s E_{p_s} = 0$.
\end{enumerate}
This representation $E$ is \new{finite-dimensional} if, for all $q \in Q_0$, $E_q$ is finite-dimensional. Its \new{dimension vector} $\vdim(E)$ is the vector $(\dim(E_q))_{q \in Q_0} \in \mathbb{N}^{\#Q_0}$. From now on, we assume that all the representations we consider are finite-dimensional. Given $E$ and $F$ two representations of $Q$, we define a \new{morphism} $\varphi : E \longrightarrow F$ as a collection $\varphi = (\varphi_q)_{q \in Q_0}$ for linear maps $\varphi_q : E_q \longrightarrow E_q$ such that, for all $\alpha \in Q_1$, we have $\varphi_{t(\alpha)} E_\alpha = F_\alpha \varphi_{s(\alpha)}$. We say that $E$ and $F$ are \new{isomorphic} whenever there exists a bijective morphism $\varphi : E \longrightarrow F$, or equivalently, if there exists a morphism $\varphi : E \longrightarrow F$ such that, for any $q \in Q_0$, $\varphi_q$ is an isomorphism of $\mathbb{K}$-vector spaces. In such a case, we write $E\cong F$. We obtain a category structure by endowing the collection of representations of $Q$ with the morphisms between them. In the following, we denote this category by $\rep(Q,R)$. In the case where $R = \varnothing$, we denote this category by $\rep(Q)$. Note that $\rep(Q, R)$ depends on the field $\mathbb{K}$.

 Any (basic) finite-dimensional $\mathbb{K}$-algebra is isomorphic to the quotient algebra $\mathbb{K}Q/\langle R \rangle$ for some bounded quiver $(Q,R)$ where $Q$ is finite, and $\langle R \rangle$ is admissible. Moreover, the category $\rep(Q,R)$ is equivalent to the finitely generated left module category over $\mathbb{K}Q/\langle R \rangle$. In particular, it implies that $\rep(Q,R)$ is an abelian Krull--Schmidt category with enough projective objects. We refer the reader to \cite{ASS06} for more details. 
 
 Write $E \oplus F$ for the direct sum of $E$ and $F$ in $\rep(Q,R)$. We said that a nonzero representation $X \in \rep(Q,R)$ is \new{indecomposable} if whenever we have $X \cong E \oplus F$ with some $E,F \in \rep(Q,R)$, then $E \cong 0$ or $F \cong 0$. Write $\ind(Q,R)$ for the collection of isomorphism classes of indecomposable representations of $(Q,R)$.

 For any $\pmb{d} = (d_q)_{q \in Q_0} \in \mathbb{N}^{\#Q_0}$, we consider the following affine space
 \[\rep((Q,R),\pmb{d}) = \prod_{\alpha \in Q_1} \Hom \left(\mathbb{K}^{d_{s(\alpha)}}, \mathbb{K}^{d_{t(\alpha)}} \right) = \prod_{\alpha \in Q_1}  \Mat_{d_{s(\alpha)} \times d_{t(\alpha)}} (\mathbb{K}),\] and refer it as the \new{representation space} of $(Q,R)$ with dimension vector $\pmb{d}$. Indeed, by choosing a basis for each of the vector spaces, we can identify the representations of $(Q,R)$ with the points of $\rep((Q,R), \pmb{d})$. The isomorphism classes of representations with vector dimension $\pmb{d}$ correspond exactly to the orbits of the group action given by the algebraic group $\vGL_{\pmb{d}}(\mathbb{K}) = \prod_{q \in Q_0} \GL_{d_q}(\mathbb{K})$, which acts on $\rep((Q,R), \pmb{d})$ by changing bases at each vertex.

\subsection{Jordan recoverability}
\label{ss:JR}

Let $E \in \rep(Q,R)$. A \new{nilpotent endomorphism} of $E$ is a morphism $N : E \longrightarrow E$  such that there exists a $k \in \mathbb{N}^*$ such that $N^k = 0$, or equivalently, for all $q \in Q_0$, the linera map $N_q$ is nilpotent. We denote the set of nilpotent endomorphisms of $E$ by $\NEnd(E)$. \cite{GPT19} shows that $\NEnd(E)$ is an irreducible algebraic variety.

An \new{integer partition} is a finite weakly decreasing sequence $\lambda = (\lambda_1,\ldots,\lambda_k)$ of positive integers. We also consider, by convention, that $(0)$ is an integer partition called \emph{zero partition}. The \new{parts} of $\lambda$ are the values, counted with multiplicities, taken by $\lambda$. The \new{size} of $\lambda$, denoted by $|\lambda|$, is the sum of all its parts; meaning $|\lambda| = \lambda_1 + \ldots + \lambda_k$. More precisely, if $|\lambda| = m$ for some $m \in \mathbb{N}$, we say that $\lambda$ is an integer partition of $m$, and we write $\lambda \vdash m$. The \new{length} of $\lambda$, denoted by $\ell(\lambda)$, is the number of parts of $\lambda$; meaning $\ell(\lambda) = k$. Let $m \in \mathbb{N}$. We define the \new{dominance order} $\leqslant$ on integer partitions of $m$ as follows: for any $\lambda, \mu \vdash m$, we say that $\lambda \leqslant \mu$ whenever for all $1 \leqslant i \leqslant \min(\ell(\lambda), \ell(\mu))$, we have $\lambda_1 + \ldots + \lambda_i \leqslant \mu_1 + \ldots + \mu_i$. 

Let $p \in \mathbb{N}^*$. We extend some notions to $p$-tuples of integer partitions, which is called \new{p-integer partitions}. Let $\pmb{\lambda} = (\lambda^1, \ldots, \lambda^p)$ be a $p$-integer partition. The \new{size vector} of $\pmb{\lambda}$ is $|\pmb{\lambda}| = (|\lambda^1|, \ldots, |\lambda^p|)$. Given $\pmb{d} = (d_1,\ldots,d_p) \in \mathbb{N}^p$, we say that $\pmb{\lambda}$ is a $p$-integer partition of $\pmb{\lambda}$ whenever $|\pmb{\lambda}| = \pmb{d}$. In such case, we write $\pmb{\lambda} \pmb{\vdash} \pmb{d}$. We define the \emph{dominance order} on $p$-integer partitions of $\pmb{d}$ as follows: for any $\pmb{\lambda},\pmb{\mu} \pmb{\vdash} \pmb{d}$, we say that $\pmb{\lambda} \domleq \pmb{\mu}$ if, for all $j \in \{1,\ldots,p\}$, we have $\lambda^j \leqslant \mu^j$.

Let $N \in \NEnd(E)$. For any $q \in Q_0$, we denote $\JF(N_q)$ the \emph{Jordan form} of the nilpotent linear map $N_q$, understood as an integer partition whose parts are the sizes of the Jordan blocks. So $\JF(N_q) \vdash d_q$ for any $q \in Q_0$. Write $\vJF(N)$ for the \new{Jordan form} of $N$, defined as the $\#Q_0$-integer partitions $(\JF(N_q))_{q \in Q_0}$. Note that $\vJF(N) \pmb{\vdash} \vdim(E)$.

\begin{theorem}[\cite{GPT19}] \label{thm:GenJFexists} Let $(Q,R)$ be a bounded quiver, $\mathbb{K}$ be an algebraically closed field, and $E \in \rep(Q,R)$. Then there exists a maximum value of $\vJF$ on $\NEnd(E)$ and it is attained in a dense open (for Zariski topology) subset of $\NEnd(E)$.
\end{theorem}

We define $\GenJF(E)$ as the maximal value of $\vJF$ on $\NEnd(E)$. We call it \new{generic Jordan form} of $E$. Note that $\GenJF$ defines an invariant on isomorphism classes of representations of $(Q,R)$. Still, it is not a complete invariant in general: that is, we can find $E,F \in \rep(Q,R)$ such that $E \ncong F$ and $\GenJF(E) = \GenJF(F)$. However, we can focus on the subcategories of $\rep(Q,R)$ such that $\GenJF$ is complete.

\begin{definition} \label{def:JR}
    Let $(Q,R)$ be a bounded quiver. A subcategory $\mathscr{C}$ or $\rep(Q,R)$ is said to be \new{Jordan recoverable} if, for any $\#Q_0$-integer partition $\pmb{\lambda}$, there exists at most one $E \in \rep(Q,R)$, up to isomorphism, such that $\GenJF(E) = \pmb{\lambda}$.
\end{definition}

To check whether a subcategory of $\rep(Q,R)$ is Jordan recoverable or not leads us to solve a system of tropical equations, which could be challenging. The following section introduces an algebraic way to build a representation from any $\#Q_0$-integer partition. Under some more restrictions, it allows one to return a representation of $(Q,R)$ from its generic Jordan form.

\subsection{Canonical Jordan recoverability}
\label{ss:CJR}

Let $\pmb{\lambda}$ be a $\#Q_0$-integer partition, and we set $|\pmb{\lambda}| = \pmb{d}$. Consider $N = \left(N_q : \mathbb{K}^{d_q} \longrightarrow \mathbb{K}^{d_q} \right)_{q \in Q_0}$ be a $\#Q_0$-tuple of nilpotent linear maps such that $\JF(N_q) = \lambda^q$. We say that a representation $E$ is \new{compatible with} $N$ whenever there exists $N' \in \NEnd(E)$ such that, for all $q \in Q_0$, the linear maps $N'_q$ and $N_q$ are similar. We denote by $\rep((Q,R),N)$ the collection of representations compatible with $N$.

We ask whether there exists a dense open set (for Zariski topology) in $\rep((Q,R),N)$ such that all its representations are isomorphic. In such a case, we define $\GenRep(N)$ as a representation in the corresponding dense open set. One can note that $\GenRep(N)$ only depends on its Jordan form. Therefore, we define $\GenRep(\pmb{\lambda})$ the \new{generic representation} of $\pmb{\lambda}$ to be one of the $\GenRep(N)$ where $N$ is a collection of nilpotent linear maps compatible with a nilpotent endomorphism $N' \in \NEnd(E)$ such that $\JF(N') = \pmb{\lambda}$.

\begin{definition} \label{def:CJR} A subcategory $\mathscr{C} \subseteq \rep(Q,R)$ is said to be \new{canonically Jordan recoverable} if, for any $X \in \mathscr{C}$, $\GenRep(\GenJF(X))$ is well-defined and is isomorphic to $X$.
\end{definition}

Note that, in general, given a $\#Q_0$-integer partition, it could happen that $\GenRep(\pmb{\lambda})$ cannot be defined. We refer to \cite[Examples 2.21 and 2.23]{D22} for examples in the case where $(Q,R)$ is a gentle quiver.

The following result shows that, at least, under some restrictions, we can ensure that $\GenRep(\pmb{\lambda})$ is well-defined for any $\#Q_0$-integer partition.

\begin{prop}[\cite{GPT19,D23}] \label{prop:GenRepDynkin}
    Let $Q$ be an $ADE$ Dynkin type quiver. Then, for any $\#Q_0$-integer partition $\pmb{\lambda}$, $\GenRep(\pmb{\lambda})$ is well-defined.
\end{prop}

In the following section, we restrict ourselves to type $A_n$ quivers for some $n \in \mathbb{N}^*$, and we recall beneficial results that reveal to be strongly linked to the $\GS$ operator on some $\mathcal{E}$-exact structures.
